% ---------------------------------------------------------------------
\section{Grammer--Yennie 分解}\label{sec:method-gy}

Grammer--Yennie 分解把胶子投影\cref{eq:method-PA}的纵横分离写成传播子分子的恒等式，并把 $i0$ 明确下来。

%%FIG fig:method-kg%%

\begin{lemma}[Grammer--Yennie 分解]\label{lem:method-gy}
(a)（共线）设胶子动量 $k$ 由 $A$ 共线子图 $C$ 流入硬部分 $H$，$\mu$ 为 $H$ 一端的指标、$\nu$ 为 $C$ 一端的指标。则
\begin{equation}\label{eq:method-gy-coll}
g_{\mu\nu}=K_{\mu\nu}+G_{\mu\nu},\qquad K_{\mu\nu}=\frac{k_\mu\,n_{B\nu}}{n_B\cdot k-i0},\qquad G_{\mu\nu}k^\nu=0,\qquad
H\cdot G\cdot C=H\cdot\Bigl(C-\frac{C^+}{k^+}\,k\Bigr) ,
\end{equation}
其中 $H\cdot G\cdot C\equiv H^\mu G_{\mu\nu}C^\nu$，$C^+=n_B\cdot C$。若 $H\cdot k=0$，则
$H\cdot C=H^+\bigl(C^--\frac{k^-}{k^+}C^+\bigr)-\mathbf H_T\cdot\bigl(\mathbf C_T-\frac{\mathbf k_T}{k^+}C^+\bigr)$。
(b)（软）设软胶子 $k$ 连接 $A$ 共线子图 $J$（指标 $\mu$）与软子图 $S$（指标 $\nu$）。则 $g^{\mu\nu}=K_s^{\mu\nu}+G_s^{\mu\nu}$，
\begin{equation}\label{eq:method-gy-soft}
\begin{gathered}
K_s^{\mu\nu}=\frac{k^\mu n_A^\nu}{n_A\cdot k\pm i0},\qquad J\cdot K_s\cdot S=\frac{(J\cdot k)\,S^-}{k^-},\\
J\cdot G_s\cdot S=J^-S^+-\mathbf J_T\cdot\mathbf S_T-\frac{(J^-k^+-\mathbf J_T\cdot\mathbf k_T)\,S^-}{k^-} ,
\end{gathered}
\end{equation}
入射的 $A$ 线取 $+i0$，出射的取 $-i0$；纯程函流 $J^\mu\propto n_A^\mu$ 被 $G_s$ 零化。
(c) \cref{eq:method-gy-coll}中的 $-i0$ 对应于 K 胶子附着在末态线上（DIS，未来指向）；附着在初态线上时（Drell--Yan）取 $+i0$。
\uses{def:method-collinear-approx,def:method-soft-approx}
\end{lemma}
% VF-WL-10, VF-WL-11

\begin{proof}
(a) $G_{\mu\nu}k^\nu=k_\mu-k_\mu(n_B\cdot k)/(n_B\cdot k)=0$（$n_B\cdot k\neq0$ 处）；$H^\mu K_{\mu\nu}C^\nu=(H\cdot k)(n_B\cdot C)/k^+=(H\cdot k)C^+/k^+$，
故 $H\cdot G\cdot C=H\cdot C-(H\cdot k)C^+/k^+=H\cdot(C-\frac{C^+}{k^+}k)$。若 $H\cdot k=0$，K 项为零，把 $H\cdot C=H^+C^-+H^-C^+-\mathbf H_T\cdot\mathbf C_T$ 中的
$H^-=(\mathbf H_T\cdot\mathbf k_T-H^+k^-)/k^+$ 代入即得所述形式。(b) 同理，$J_\mu K_s^{\mu\nu}S_\nu=(J\cdot k)(n_A\cdot S)/(n_A\cdot k)=(J\cdot k)S^-/k^-$；
$J\cdot G_s\cdot S=J\cdot S-(J\cdot k)S^-/k^-$，把 $J\cdot S=J^+S^-+J^-S^++\ldots$ 与 $J\cdot k=J^+k^-+J^-k^+-\mathbf J_T\cdot\mathbf k_T$ 代入，$J^+S^-$ 项相消。
若 $J\propto n_A$，则 $J^-=\mathbf J_T=0$，右边为零。(c) 见\cref{sec:appF-gy}：$i0$ 由被 K 项取代的传播子的极点位置决定。这些恒等式都用符号计算核对过。% check: w10_gy.py
\end{proof}

最初的 QED 形式（\cite{GrammerYennie1973}，如\cite{SenIndumathiChoudhury2019}（预印本）的式~(7)--(9)所转述）中 K 部分正比于 $k_\mu k_\nu$，系数由相邻带电线的传播子构成，全部红外发散都在 K 光子中；
带固定辅助矢量与指定 $i0$ 的单侧形式是 Collins 的推广（\cite[第~5.5.6、10.6、11.9.3~节]{Collins2011}），\cref{eq:method-gy-coll}即\cite[式~(11.4)--(11.6)，第~402~页]{Collins2011}，
须按作者的勘误表交换上下指标（附录~Q 第~19~条，\cref{tab:appQ-errata-b}）；历史与 $i0$ 的讨论见\cref{sec:appF-gy}。


\begin{proposition}[G 部分的幂次压低 \tagA]\label{prop:method-gy-suppression}
（\cref{fig:method-kg}）在共线近似之后：除最小部分子内容之外，每多一条共线 G 胶子至少压低一个 $\lambda$；K 胶子不带来任何压低，任意多条都是领头的。
唯一的例外是作为活动部分子的单个横向胶子：此时 $k\cdot H=0$，K 项为零而 G 项是领头的，并经由组合 $k^+A^j-k^jA^+$ 组成胶子分布中的场强 $G^{+j}$（\cref{prop:ir-super-leading}）。
软 G 胶子在 Glauber 区之外压低 $O(\lambda)$，在 Glauber 区中为 $O(1)$。（本命题只涉及部分子态。）
% errtype: E-S
\uses{lem:method-gy,prop:ir-power-counting,prop:ir-super-leading}
\end{proposition}

\begin{proofsketch}
(1) 由推动论证（\cref{prop:ir-power-counting}），$A$ 共线子图的流按 $C^\nu\sim(1,\lambda^2,\lambda)\,C^+$ 标度，而 $k\sim(1,\lambda^2,\lambda)\,k^+$；
于是 $C-\frac{C^+}{k^+}k=(0,O(\lambda^2),O(\lambda))\,C^+$，由\cref{lem:method-gy}(a)，$H\cdot G\cdot C=O(\lambda)\,H\cdot C$，而 K 项 $H\cdot K\cdot C=(H\cdot k)C^+/k^+$ 不含压低因子。
(2) 若唯一的共线部分子是胶子，Ward 恒等式给出 $k\cdot H=0$（其余外线物理），\cref{lem:method-gy}(a) 的第二式中领头的是 $-\mathbf H_T\cdot(\mathbf C_T-\frac{\mathbf k_T}{k^+}C^+)=-H^j(k^+C^j-k^jC^+)/k^+$，
即场强组合 $k^+A^j-k^jA^+$；这与\cref{prop:ir-super-leading}中超领头构型的相消一致。(3) 软胶子：由\cref{eq:method-gy-soft}，$J\cdot G_s\cdot S$ 相对于 $J\cdot K_s\cdot S$ 为 $O(\lambda|\mathbf k_T|/|k^-|)$，
软区为 $O(\lambda)$，Glauber 区为 $O(1)$。(4) 数值上（\cref{sec:appF-gy}）：单圈 DIS 顶点中，在壳无质量母夸克配物理光子时 G 项按 $\lambda^2$ 压低（拟合幂次 $2.00$），一般情形按 $\lambda$ 压低（拟合幂次 $0.99$--$1.00$）；
软区的 $G_s$ 项按 $\lambda$ 压低（拟合幂次 $1.00$），Glauber 区出现 $O(1)$ 的平台。（以上只给出 E-S 型陈述。）
\end{proofsketch}

\paragraph{地位与缺口。}全部是标度陈述（E-S）。积分形式需要\cref{asm:appN-envelope}（\tagA）：误差表第 3、4、6 行正是“相对于包络”的投影误差，(GA3)(iii) 是推动论证的逐点形式，
而推动论证本身只给出标度。胶子道的例外说明“横向胶子的数目”不是规范不变的计数（\cref{prop:ir-super-leading}）。文献：\cite[第~11.3~节]{Collins2011}；
Collins--Soper--Sterman\cite[第~8、9.1~节]{CollinsSoperSterman1989}指出，把硬部分的外动量换成类光动量并只保留领头的极化分量之后，连接硬部分与喷注的额外共线胶子恰好是纵向极化的，修正被 $Q$ 的幂次压低——这正是 K 部分领头、G 部分压低的标度陈述。
